\documentclass[Screen16to9,17pt]{foils}
\usepackage{zencurity-slides}
\externaldocument{software-security-exercises}
\selectlanguage{english}

\begin{document}

\mytitlepage
{13. Security Design}
{EK  Software Security}

\slide{Plan for today}

\begin{list1}
\item Subjects
\begin{list2}
\item How to reach the goal of secure design
\item Security Principles for software
\item Principle of least privilege, fail-safe defaults, separation of privilege etc.
\item Files, objects, users, groups and roles
\item Security by Design
\item Privacy by Design
\item Benchmarking standards
\item CIS controls Center for Internet Security
\item Cases and your input!
\end{list2}
\item Exercises
\begin{list2}
\item How should software be designed today
\end{list2}
\end{list1}

\slide{Time schedule}

\hlkimage{4cm}{kame-noanime-small.png}

Plan for the day
\begin{list1}
\item 12:45 - 13:30 -- Session 0: Security Design
\item 13:30 - 13:45 -- 15min break
\item 13:45 - 14:45 -- Session 1+2: Microsoft Windows and Apple Security
\item 14:45 - 15:00 -- 15min break
\item 15:00 - 16:45 -- Session 3: Working with CIS Benchmarking
\end{list1}


\slide{Session 0: Security Design}

I would like for you to read some more documents ...


\begin{list1}
\item We will use these as a reference
\item Security by Design Principles\\
\url{https://www.owasp.org/index.php/Security_by_Design_Principles}
\item Selected as OWASP has existed for many years, expect site to survive

\item Likewise we will use the Wikipedia article about \emph{Privacy by Design}\\
\url{https://en.wikipedia.org/wiki/Privacy_by_design}
\end{list1}


\slide{Other resources used today}

\begin{list2}
\item ENISA, EU security office, \url{https://www.enisa.europa.eu/publications/privacy-and-data-protection-by-design}\\ and \url{https://www.enisa.europa.eu/publications/big-data-protection}
\item  \emph{Windows 11 Security Book: Powerful security by design}, Microsoft \\
Available in the folder as: \verb+MSFT-Windows11-Security-book_Sept2023.pdf+
\item \emph{Apple Platform Security}, august 2026, 266 pages, Apple\\
\url{https://help.apple.com/pdf/security/en_US/apple-platform-security-guide.pdf}
\end{list2}


\vskip 2cm
\begin{center}
On purpose using references that have existed for many years. Secure design has not changed much, but developers have -- they are sometimes \emph{lazy}
\end{center}

\slide{Goals: Foundation and design}

.
\hlkrightimage{8cm}{kyler-trautner-693525-unsplash.jpg}

\begin{list2}
\item How to reach the goal of secure design
by listing Security Principles for \\
software like Principle of least privilege, fail-safe defaults, separation of\\
privilege etc.
\item Security by Design and Privacy by Design requires a lot of thought and\\ perhaps Benchmarking standards
like CIS controls Center for Internet\\
 Security can help
\item Security Principles for software
\item Some repetition, hope it gives a nice roundup for the course
\end{list2}

\hfill {\footnotesize Picture from Unsplash kyler-trautner-693525-unsplash.jpg}

\slide{Assumptions}

\begin{quote}
Any security policy, mechanism, or procedure is based on assumptions that, if incorrect, destroy the superstructure on which it is built.\\
Matt Bishop, Computer Security 2019
\end{quote}

\begin{list1}
\item Example, vendor patches
\item Important points:
\begin{list2}
\item Is the patch correct? Example Spectre and Heartbleed
\item Vendor test environments equal to intended environments
\item Installed correctly - including operator skills
\end{list2}
\end{list1}


\slide{Weak Structural Security}

Example design flaws:
\begin{list2}
\item Large Attack surface
\item Running a Process at Too High a Privilege Level, dont run everything as root or administrator
\item No Defense in Depth, use more controls, make a strong chain
\item Not Failing Securely
\item Mixing Code and Data
\item Misplaced trust in External Systems
\item Insecure Defaults
\item Missing Audit Logs
\end{list2}

\slide{Secure Programming for Linux and Unix Howto}

\begin{list1}
\item More information about systems design and implementation can be found in the free resource:
\item Secure Programming for Linux and Unix HOWTO, David Wheeler
\item \link{https://dwheeler.com/secure-programs/Secure-Programs-HOWTO.pdf}
\item Chapter 5. Validate All Input details input validation in the context of Unix programs
\item Chapter 6. Restrict Operations to Buffer Bounds (Avoid Buffer Overflow)
\item Chapter 7. Design Your Program for Security
\end{list1}

\slide{Example applications from Microsoft}

How to get ahead? - use existing good examples!

Microsoft has released sample applications.

\begin{quote}
Secure Development Documentation
Learn how to develop and deploy secure applications on Azure with our sample apps, best practices, and guidance.

Get started
Develop a secure web application on Azure
\end{quote}

Source:
\url{https://docs.microsoft.com/en-us/azure/security/develop/}

Yes, this describes how to run Alpine Linux on their Azure Cloud.


\slide{Defense in depth}

%\hlkimage{10cm}{Bartizan.png}
\hlkimage{14cm}{medieval-clipart-5}
\centerline{Picture originally from: \url{http://karenswhimsy.com/public-domain-images}}


\slide{Defense in depth - layered security}

\hlkimage{5cm}{security-layers-1-uk.pdf}

\centerline{\hlkbig\color{security6blue} Multiple layers of security! Isolation!}



\slide{Security Principles}

We will use these as a reference:
\begin{list2}
\item The Protection of Information in Computer Systems is a 1975 seminal publication by Jerome Saltzer and Michael Schroeder about information security. The paper emphasized that the primary concern of security measures should be the information on computers and not the computers themselves.\\
\url{https://en.wikipedia.org/wiki/The_Protection_of_Information_in_Computer_Systems}
\item The paper lists eight design principles: economy of mechanism, fail-safe defaults, complete mediation, open design, separation of privilege, least privilege, least common mechanism and psychological acceptability
\item Principle of least privilege -- individual principles often have their own pages:\\
\url{https://en.wikipedia.org/wiki/Principle_of_least_privilege}
\item See also \emph{Security by Design Principles}\\
Selected as OWASP has existed for many years, expect site to survive\\
\url{https://www.owasp.org/index.php/Security_by_Design_Principles}
\end{list2}


\slide{Principle of Least Privilege}

\begin{list1}
\item {\bf Definition 14-1} The \emph{principle of least privilege} states that a subject should be given only those privileges that it needs in order to complete the task.
\item Also drop privileges when not needed anymore, relinquish rights immediately
\item Example, need to read a document - but not write.
\item Database systems can often provide very fine grained access to data
\end{list1}

{\footnotesize Source: \emph{Computer Security}, Matt Bishop, 2019}


\slide{Principle of Fail-Safe defaults}

\begin{list1}
\item {\bf Definition 14-3} The \emph{principle of fail-safe defaults} states that, unless a subject is given explicit access to an object, it should be denied access to that object.
\item Default access \emph{none}
\item In firewalls default-deny - that which is not allowed is prohibited
\item Newer devices today can come with no administrative users, while older devices often came with default admin/admin users
\item Real world example, OpenSSH config files that come with \verb+PermitRootLogin no+
\end{list1}

{\footnotesize Source: \emph{Computer Security}, Matt Bishop, 2019}


\slide{Principle of Economy of Mechanism}

\begin{list1}
\item {\bf Definition 14-4} The \emph{principle of economy of mechanism} states that security mechanisms should be as simple as possible.
\item Simple $->$ fewer complications $->$ fewer security errors
\item Use WPA passphrase instead of MAC address based authentication
\item
\end{list1}

{\footnotesize Source: \emph{Computer Security}, Matt Bishop, 2019}


\slide{Principle of Complete Mediation}

\begin{list1}
\item {\bf Definition 14-5} The \emph{principle of complete mediation} requires that all accesses to objects be checked to ensure that they are allowed.
\item Always perform check
\item Time of check, time of use
\item Example Unix file descriptors - access check first, then can be reused in the future
\item Caching can be bad.
\end{list1}

{\footnotesize Source: \emph{Computer Security}, Matt Bishop, 2019}


\slide{Principle of Open Design}

\hlkimage{6cm}{debath-stego.png}

Source: picture from \link{https://www.cs.cmu.edu/~dst/DeCSS/Gallery/Stego/index.html}
\begin{list1}
\item {\bf Definition 14-6} The \emph{principle of open design} states that the security of a mechanism should not depend on the secrecy of its design or implementation.
\item Content Scrambling System (CSS) used on DVD movies
\item Mobile data encryption  A5/1 key
\end{list1}

{\footnotesize Source: \emph{Computer Security}, Matt Bishop, 2019}


\slide{Security by Design}

\begin{list1}
\item All the principles today lead to {\bfseries Security by Design}
\begin{list2}
\item Secure defaults, minimize attack surface, fail securely
\item Least privilege -- no more ``everything runs as root''
\item Security decisions are explicit choices, not accidents
\end{list2}
\item Applications don't run in a vacuum
\begin{list2}
\item They run in existing environments -- integrate with existing solutions, don't expect a clean slate
\item Make it easy to use by following conventions -- comply or explain why not
\end{list2}
\item Security from requirements and design to deployment
\begin{list2}
\item Threat modeling during design, checklists and cheat sheets to avoid repeating mistakes
\end{list2}
\end{list1}

Source: OWASP Developer Guide, 2.3 Principles of security and 4 Design\\
{\footnotesize\link{https://owasp.org/www-project-developer-guide/release/}}\\
See also: ENISA Secure by Design and Default Playbook


\slide{Breaking News: ENISA Secure by Design and Default Playbook}

{\bfseries\Large Hot off the press!} -- well, almost: published 30 July 2026

\begin{list1}
\item ENISA Secure by Design and Default Playbook\\
A Practical Guide to Secure by Design and Default Principles for SMEs
\begin{list2}
\item Principles turned into clear, repeatable actions
\item Fits into existing engineering, product and release processes
\item Covers the whole product life cycle
\item 22 playbooks, browsable on GitHub
\end{list2}
\item {\footnotesize\link{https://www.enisa.europa.eu/publications/enisa-secure-by-design-and-default-playbook}}
\item {\footnotesize\link{https://github.com/enisaeu/enisa-sbd-playbook/}}
\end{list1}

Today's principles are not old news -- the EU is still writing about them!


\slide{Files, objects, users, groups and roles}

\begin{quote}
Establish secure defaults
There are many ways to deliver an “out of the box” experience for users. However, by default, the experience should be secure, and it should be up to the user to reduce their security – if they are allowed.

For example, by default, password aging and complexity should be enabled. Users might be allowed to turn these two features off to simplify their use of the application and increase their risk.
\end{quote}

Source: OWASP Security by Design Principles\\
\url{https://www.owasp.org/index.php/Security_by_Design_Principles}

\begin{list2}
\item Applications dont run in vacuum, they run in existing environments
\item Make sure to integrate with existing solutions, dont expect a clean slate
\item Make it easy to use by following convention
\end{list2}

\slide{Fail securely}

\begin{minted}[fontsize=\footnotesize]{java}
isAdmin = true;
try {
  codeWhichMayFail();
  isAdmin = isUserInRole( “Administrator” );
}
catch (Exception ex) {
  log.write(ex.toString());
}
\end{minted}
"If either codeWhichMayFail() or isUserInRole fails or throws an exception, the user is an admin by default. This is obviously a security risk."\\
Source: OWASP Security by Design Principles\\
\url{https://www.owasp.org/index.php/Security_by_Design_Principles}

\begin{list2}
\item DO use exception handling!
\item Just make sure code does not end up giving more privileges when failing!
\end{list2}

\slide{Separation of Duties and Function}

\begin{quote}
{\bf Separation of duties} (SoD; also known as Segregation of Duties) is the concept of having more than one person required to complete a task. In business the separation by sharing of more than one individual in one single task is an internal control intended to prevent fraud and error.
\end{quote}

Quote from \url{https://en.wikipedia.org/wiki/Separation_of_duties}

\begin{quote}
{\bf Separation of function}. Developers do not develop new programs on production systems because of the potential threat to production data.
\end{quote}
\emph{Computer Security}, Matt Bishop, 2019

Danish: Funktionsadskillelse


\slide{Role-based Access Control (RBAC)}

\begin{quote}
In computer systems security, {\bf role-based access control (RBAC)}[1][2] or role-based security[3] is an approach to restricting system access to unauthorized users. It is used by the majority of enterprises with more than 500 employees,[4] and can implement mandatory access control (MAC) or discretionary access control (DAC).

Role-based access control (RBAC) is a policy-neutral access-control mechanism defined around {\bf roles and privileges}. The components of RBAC such as role-permissions, user-role and role-role relationships make it simple to perform user assignments. A study by NIST has demonstrated that RBAC addresses many needs of commercial and government organizations[citation needed]. RBAC can be used to facilitate administration of security in large organizations with hundreds of users and thousands of permissions. Although RBAC is different from MAC and DAC access control frameworks, it can enforce these policies without any complication.
\end{quote}
Quote from \url{https://en.wikipedia.org/wiki/Role-based_access_control}



\slide{Example role based system: PostgreSQL security}
\hlkimage{15cm}{postgresql-security.png}

Feature overview security features in PostgreSQL\\
\url{https://www.postgresql.org/about/featurematrix/#security}


\slide{Security by Design}

\hlkimage{6cm}{OWASP-SDLC.jpg}

\begin{list2}
\item Doing the above should result in applications which are secure by design
\item Adhering to the best security principles
\item Implementing security from design to deployment ensure good security
\end{list2}


\slide{Privacy by Design}

\begin{quote}
Objectives of the report\\
This report shall promote the discussion on how privacy by design can be implemented with the help
of engineering methods. It provides a basis for better understanding of the current state of the art
concerning privacy by design with a focus on the technological side.

“Personal data” means any information relating to an identified or identifiable natural person—for a
detailed discussion see [19]. This is related to the term personally identifiable information (PII), as e.g.
used in the privacy framework standardised by ISO/IEC [125].
\end{quote}

Source: Privacy and Data Protection by Design -- from policy to engineering\\
{\footnotesize\url{https://www.enisa.europa.eu/publications/privacy-and-data-protection-by-design}}

\begin{list2}
\item Next, if the application is secure, what about handling and protecting personal data
\end{list2}

\slide{ENISA: List of Recommendations}

\begin{list2}
\item {\small Policy makers need to support the development of new incentive mechanisms for privacy-friendly services and need to promote them.}
\item {\small The research community needs to further investigate in privacy engineering, especially with a
multidisciplinary approach. This process should be supported by research funding agencies.
The results of research need to be promoted by policy makers and media.}
\item {\small Providers of software development tools and the research community need to offer tools that
enable the intuitive implementation of privacy properties.}
\item {\small Especially in publicly co-founded infrastructure projects, privacy-supporting components,
such as key servers and anonymising relays, should be included.}
\item {\small Data protection authorities should play an important role providing independent guidance
and assessing modules and tools for privacy engineering.}
\item {\small Legislators need to promote privacy and data protection in their norms.}
\item {\small Standardisation bodies need to include privacy considerations in the standardisation process.}
\item {\small Standards for interoperability of privacy features should be provided by standardization bodies.}
\end{list2}

Source: Privacy and Data Protection by Design -- from policy to engineering\

\slide{Privacy Protection Goals}

\begin{quote}\footnotesize
In ICT
security the triad of confidentiality, integrity, and availability has been widely accepted. ...
As complement to these security protection goals, three privacy-specific protection goals were pro-
posed in 2009 [172], namely unlinkability, transparency, and intervenability.
\end{quote}

\begin{list2}
\item {\small \emph{Unlinkability}. Unlinkability ensures that privacy-relevant data cannot be linked across domains that are constituted by a common purpose and context, and that means that processes have to be oper-
ated in such a way that the privacy-relevant data are unlinkable to any other set of privacy-
relevant data outside of the domain.}
\item {\small \emph{Transparency.} Transparency ensures that all privacy-relevant data processing including the
legal, technical and organisational setting can be understood and reconstructed at any time.
The information has to be available before, during and after the processing takes place. Thus,
transparency has to cover not only the actual processing, but also the planned processing (ex-
ante transparency) and the time after the processing has taken place to know what exactly
happened (ex-post transparency).}
\item {\small \emph{Intervenability.} Intervenability ensures intervention is possible concerning all ongoing or
planned privacy-relevant data processing, in particular by those persons whose data are pro-
cessed. The objective of intervenability is the application of corrective measures and counter-
balances where necessary.}
\end{list2}
Source: Privacy and Data Protection by Design -- from policy to engineering



\slide{Data oriented strategies}

\begin{list2}
\item Strategy \#1: MINIMISE\\
The most basic privacy design strategy is MINIMISE, which states that the amount of personal data
that is processed should be restricted to the minimal amount possible.
\item Strategy \#2: HIDE\\
The second design strategy, HIDE, states that any personal data, and their interrelationships, should
be hidden from plain view.
\item Strategy \#3: SEPARATE\\
The third design strategy, SEPARATE, states that personal data should be processed in a distributed
fashion, in separate compartments whenever possible.
\item Strategy \#4: AGGREGATE\\
The fourth design pattern, AGGREGATE, states that Personal data should be processed at the highest
level of aggregation and with the least possible detail in which it is (still) useful.
\end{list2}

Source: Privacy and Data Protection by Design -- from policy to engineering



\slide{Privacy Impact Assessment (PIA)}

\begin{quote}\small
Privacy-by-design er en tilgang, der sikrer, at virksomheden indarbejder databeskyttelse som en integreret del af virksomhedens forretningsprocesser, værdikæde og produktlivscyklus. Lige fra produktionsfasen til at produktet havner hos slutbrugeren.
...
Udarbejd en konsekvensanalyse (PIA)\\
En god tilgang til at sikre privacy-by-design/default er, at udarbejde en konsekvensanalyse også kaldet en PIA (Privacy Impact Assessment), der tydeliggør konsekvenserne af virksomhedens databehandling.
\end{quote}

En PIA skal som minimum indeholde følgende:

\begin{list2}
\item En beskrivelse af de påtænkte databehandlinger og deres formål.
\item En vurdering af behandlingens nødvendighed og proportionalitet.
\item En vurdering af risikoen for de personer, hvis persondata bliver behandlet.
\item Foranstaltninger til at imødegå risici og demonstrere overholdelse af persondataforordningen.
\end{list2}

Source: IT-Branchen: Privacy-by-design/default\\ \url{https://itb.dk/persondataforordningen/privacy-by-design-default/}


\slide{Security Controls and Frameworks}

\begin{list1}
\item Multiple exist
\vskip 1cm
\begin{list2}
\item CIS controls Center for Internet Security (CIS) \link{https://www.cisecurity.org}
\item PCI Best Practices for Maintaining PCI DSS Compliance v2.0 Jan 2019
\item NIST Cybersecurity Framework (CSF)\\
Framework for Improving
Critical Infrastructure Cybersecurity\\ \link{https://www.nist.gov/cyberframework}\\
\link{http://csrc.nist.gov/publications/PubsSPs.html}
\item National Security Agency (NSA)\\ \link{http://www.nsa.gov/research/publications/index.shtml}
\item Information Systems Audit and Control Association (ISACA)\\
\link{http://www.isaca.org/Knowledge-Center/Risk-IT-IT-Risk-Management/Pages/default.aspx}
\end{list2}
\end{list1}


\slide{ISO/IEC 27001:2022 Information Security Management Systems}

\begin{list2}
\item ISO/IEC 27001:2022 specifies requirements for an Information Security Management System (ISMS)
\item Annex A lists 93 reference controls, grouped in four themes:\\
5 Organizational (37), 6 People (8), 7 Physical (14), 8 Technological (34)
\item ISO/IEC 27002:2022 gives implementation guidance for each control
\item The organization decides which controls apply, based on risk assessment\\
-- documented in the Statement of Applicability (SoA)
\item Excluding a control requires a justification\\
``We don't develop software'' is often not a valid reason -- scripts, automation and configuration count too
\end{list2}

Source: ISO/IEC 27001:2022, ISO/IEC 27002:2022


\slide{Annex A -- Secure Development Controls}

\begin{list2}
\item 8.25 Secure development life cycle
\item 8.26 Application security requirements
\item 8.27 Secure system architecture and engineering principles
\item 8.28 Secure coding
\item 8.29 Security testing in development and acceptance
\item 8.30 Outsourced development
\vskip 1em
\item Your ISMS project has to answer for each: does it apply, and how do we meet it?
\end{list2}

Source: ISO/IEC 27001:2022 Annex A


\slide{Risk management defined}

\hlkimage{20cm}{shon-harris-risk-management.png}

Source: Shon Harris \emph{CISSP All-in-One Exam Guide}


\slide{Center for Internet Security CIS Controls}

\begin{quote}
  The CIS ControlsTM are a prioritized set of actions that collectively form a defense-in-depth set
of best practices that mitigate the most common attacks against systems and networks. The
CIS Controls are developed by a community of IT experts who apply their first-hand experience
as cyber defenders to create these globally accepted security best practices. The experts who
develop the CIS Controls come from a wide range of sectors including retail, manufacturing,
healthcare, education, government, defense, and others.
\end{quote}

Source: \link{https://www.cisecurity.org/} CIS-Controls-Version-7-1.pdf

\slide{CIS Controls 7.1: Five Critical Tenets}

\begin{list2}
\item {\bf Offense informs defense:} Use knowledge of actual attacks that have
compromised systems to provide the foundation to continually learn
from these events to build effective, practical defenses. Include only
those controls that can be shown to stop known real-world attacks.
\item {\bf Prioritization:} Invest first in Controls that will provide the greatest risk
reduction and protection against the most dangerous threat actors
and that can be feasibly implemented in your computing environment.
\item {\bf Measurements and Metrics:} Establish common metrics to provide a
shared language for executives, IT specialists, auditors, and security
officials to measure the effectiveness of security measures within
an organization so that required adjustments can be identified and
implemented quickly.
\item {\bf Continuous diagnostics and mitigation:} Carry out continuous
measurement to test and validate the effectiveness of current security
measures and to help drive the priority of next steps.
\item {\bf Automation:} Automate defenses so that organizations can achieve
reliable, scalable, and continuous measurements of their adherence to
the Controls and related metrics.
\end{list2}

{\footnotesize Source: Center for Internet Security CIS Controls 7.1 \verb+CIS-Controls-Version-7-1.pdf+}


\slide{CIS Controls 7.1 versus v8}

{\footnotesize
\begin{tabular}{ll}
v7.1 & v8 \\ \hline
1 Hardware Assets & 1 Enterprise Assets \\
2 Software Assets & 2 Software Assets \\
3 Continuous Vulnerability Management & 7 Continuous Vulnerability Management \\
4 Administrative Privileges & 5 Account Management, 6 Access Control Management \\
5 Secure Configuration & 4 Secure Configuration \\
6 Audit Logs & 8 Audit Log Management \\
18 Application Software Security & 16 Application Software Security \\
\end{tabular}
}

\begin{list2}
\item v7.1 has 20 controls, v8 and v8.1 have 18 controls and 153 Safeguards
\item v8 groups by activity, not by who manages the device: enterprise assets cover laptops, servers, IoT, cloud
\item New in v8: Control 3 Data Protection, Control 15 Service Provider Management
\item Prioritize with Implementation Groups: IG1 is \emph{essential cyber hygiene}, every enterprise should start with IG1
\end{list2}

{\footnotesize Source: \link{https://www.cisecurity.org/controls/implementation-groups}}


\slide{Inventory and Control of Enterprise Assets}

\begin{quote}
CIS Control 1:\\
Inventory and Control of Enterprise Assets\\
Actively manage (inventory, track, and correct) all enterprise assets [...] to accurately know the totality of assets that need to be monitored and protected within the enterprise. This will also support identifying unauthorized and unmanaged assets to remove or remediate.
\end{quote}

\begin{list1}
\item What is connected to our networks -- and running in our cloud accounts?
\item What firmware do we need to install on hardware?
\item Where IS the hardware we own?
\item What hardware is still supported by vendor?
\end{list1}

Source: CIS Controls v8 \verb+CIS_Controls_v8_Guide.pdf+, was Control 1 Hardware Assets in v7.1


\slide{Inventory and Control of Software Assets}

\begin{quote}
CIS Control 2:\\
Inventory and Control of Software Assets\\
Actively manage (inventory, track, and correct) all software (operating systems and applications) on the network so that only authorized software is installed and can execute, and that unauthorized and unmanaged software is found and prevented from installation or execution.
\end{quote}

\begin{list1}
\item What licenses do we have? Paying too much?
\item What versions of software do we depend on?
\item Which libraries are inside our own software? Generate an SBOM
\item What software needs to be phased out, upgraded?
\end{list1}

Source: CIS Controls v8 \verb+CIS_Controls_v8_Guide.pdf+, same number in v7.1


\slide{Secure Configuration of Enterprise Assets and Software}

\begin{quote}
CIS Control 4:\\
Secure Configuration of Enterprise Assets and Software\\
Establish and maintain the secure configuration of enterprise assets (end-user devices, including portable and mobile; network devices; non-computing/IoT devices; and servers) and software (operating systems and applications).
\end{quote}

\begin{list1}
\item Start from a CIS Benchmark: OS, Kubernetes, database, browser
\item Default passwords changed, unused services and debug features off?
\item Check automatically: kube-bench, OpenSCAP, Lynis
\item Keep configuration in git, review changes
\end{list1}

Source: CIS Controls v8 \verb+CIS_Controls_v8_Guide.pdf+, was Control 5 in v7.1


\slide{Account Management and Access Control Management}

\begin{quote}
CIS Control 5: Account Management\\
Use processes and tools to assign and manage authorization to credentials for user accounts, including administrator accounts, as well as service accounts, to enterprise assets and software.
\end{quote}

\begin{quote}
CIS Control 6: Access Control Management\\
Use processes and tools to create, assign, manage, and revoke access credentials and privileges for user, administrator, and service accounts for enterprise assets and software.
\end{quote}

\begin{list1}
\item Separate admin accounts from daily accounts, who owns each service account?
\item Revoke access when people leave or change role
\end{list1}

Source: CIS Controls v8 \verb+CIS_Controls_v8_Guide.pdf+, was Control 4 Administrative Privileges in v7.1


\slide{Continuous Vulnerability Management}

\begin{quote}
CIS Control 7:\\
Continuous Vulnerability Management\\
Develop a plan to continuously assess and track vulnerabilities on all enterprise assets within the enterprise's infrastructure, in order to remediate, and minimize, the window of opportunity for attackers. Monitor public and private industry sources for new threat and vulnerability information.
\end{quote}

\begin{list1}
\item Scan OS packages, container images and dependencies, e.g. \verb+trivy image+
\item Subscribe to security advisories for the software you depend on
\item Agree on fix deadlines per severity, re-scan to verify the fix
\end{list1}

Source: CIS Controls v8 \verb+CIS_Controls_v8_Guide.pdf+, was Control 3 in v7.1


\slide{Audit Log Management}

\begin{quote}
CIS Control 8:\\
Audit Log Management\\
Collect, alert, review, and retain audit logs of events that could help detect, understand, or recover from an attack.
\end{quote}

\begin{list1}
\item ... and present it, use it daily, report it to management!
\item Which events: logins, admin actions, configuration changes
\item Collect centrally, synchronize time with NTP
\item Decide how long logs are retained
\end{list1}

Source: CIS Controls v8 \verb+CIS_Controls_v8_Guide.pdf+, was Control 6 in v7.1


\slide{Application Software Security}

\begin{quote}
CIS Control 18:\\
Application Software Security\\
Manage the security life cycle of all in-house developed and acquired software in order to prevent, detect, and correct security weaknesses.
\end{quote}

Source: Center for Internet Security CIS Controls 7.1 \verb+CIS-Controls-Version-7-1.pdf+

\begin{quote}
CIS Control 16:\\
Application Software Security\\
Manage the security life cycle of in-house developed, hosted,
or acquired software to prevent, detect, and remediate security
weaknesses before they can impact the enterprise.
\end{quote}

Source: Center for Internet Security CIS Controls v8 \verb+CIS_Controls_v8_Guide.pdf+


\slide{Session 2+3: Microsoft Windows and Apple Security}

I would like for you to read some more documents ...

We will go through some documents from the resources in the folder Supporting resources online, hopefully inspire you to read more of them.

\begin{list1}
\item  \emph{Windows 11 Security Book: Powerful security by design}, Microsoft \\
Available in the folder as: \verb+MSFT-Windows11-Security-book_Sept2023.pdf+
\item \emph{Apple Platform Security}, august 2026, 266 pages, Apple\\
\url{https://help.apple.com/pdf/security/en_US/apple-platform-security-guide.pdf}
\end{list1}

Maybe not read every word, but look through the Windows doc and browser the Apple one!


\slide{Session 3: Working with CIS Benchmarking}

\begin{list2}
\item Even though the OS is updated,
\item even though the containers are updated,
\item What are secure settings
\end{list2}

How do we check?

Find and run CIS benchmark

Let's run kubebench


\slide{Kubernetes Security is more than just the core}

\hlkimage{12cm}{rice-container-security-1.png}
Source: \emph{Container Security}, (CS) Liz Rice, O'Reilly, Apr 2020

\begin{list2}
\item Attacks from inside, attacks inside containers
\item Some just use up resources, some attack others, some attack the system from the inside
\item \emph{Hacking Kubernetes} has a whole \emph{Chapter 3 Container Runtime Isolation}
\end{list2}


\slide{Exercise: browse CVE sites, find vulnerabilities in software you use}

Try visiting the sites:
\begin{list2}
\item CVE Main site \link{https://www.cve.org}
\item Nice searching and graphs \link{https://www.cvedetails.com}
\item Even more graphs and insights \link{https://zerodayclock.com/}
\item CISA Known Exploited Vulnerabilities (KEV) catalog, actively exploited\\
\link{https://www.cisa.gov/known-exploited-vulnerabilities-catalog}
\end{list2}

Please spend a little time, try searching for some vendor you know

Bookmark the sites for future reference.


\slide{Benchmarking tools}

\begin{quote}
\faWrench\ \verb+Kube-bench+ is the industry-standard tool to automate checking Kubernetes compliance with the Center for Internet Security (CIS) Benchmark.

Kube-bench makes it easy for operators to check whether each node in their Kubernetes cluster is configured according to security best practices.
\end{quote}
Source: \link{https://info.aquasec.com/open-source}

\begin{list2}
\item CIS Kubernetes Benchmark v2.1.0 - Nov 2026, Kubernetes 1.36 and 1.37 -- other versions exist
\item CIS Docker Benchmark v1.9.0 - Nov 2026, Docker v29.x
\end{list2}

%\slide{Benchmarking and keeping up to date}

% Try it out \faWrench\ \verb+kube-bench+ - make some bad changes, like LKS p95 allow anonymous authentication, show why layered defense works, and how kube-bench marks it. Another example token based access, initially we have token while installing, remember to remove!
% \faWrench\ \verb+kube-hunter+ and kube-bench

\slide{Tool example kube-bench}

%\hlkimage{}{}

\begin{alltt}\scriptsize
hlk@timon:~/bin/kube-bench/kube-bench$ {\bf kubectl logs kube-bench-gdf62}
[PASS] 1.1.7 Ensure that the etcd pod specification file permissions are set to 600 or more restrictive (Automated)
[PASS] 1.1.8 Ensure that the etcd pod specification file ownership is set to root:root (Automated)
[WARN] 1.1.9 Ensure that the Container Network Interface file permissions are set to 600 or more restrictive (Manual)
[WARN] 1.1.10 Ensure that the Container Network Interface file ownership is set to root:root (Manual)
[PASS] 1.1.11 Ensure that the etcd data directory permissions are set to 700 or more restrictive (Automated)
[FAIL] 1.1.12 Ensure that the etcd data directory ownership is set to etcd:etcd (Automated)
...
== Summary policies ==
0 checks PASS
0 checks FAIL
35 checks WARN
0 checks INFO

== Summary total ==
63 checks PASS
10 checks FAIL
58 checks WARN
0 checks INFO

\end{alltt}

\begin{list2}
\item \url{https://github.com/aquasecurity/kube-bench} also check out Lynis \url{https://cisofy.com/lynis/}
\end{list2}


\exercise{ex:sw-kubebench-k3s}

\slidenext{}


\end{document}
